\chapter{Engineering Organisation Design}\label{ch:fm:engineering-organisation-design}

In April 1968 Melvin Conway published the observation that ``organizations which design systems \dots\ are constrained to produce designs which are
copies of the communication structures of these organizations''. A trading firm that splits its engineers by asset class ends up with three order
gateways, three risk checks and three market-data handlers, whether or not it wanted them, because three teams that rarely talk build three of
everything. The organisation chart is the first draft of the architecture.

\section{Embedded against central teams}

\begin{definition}[Embedded team, central team]\label{def:fm:engineering-organisation-design:teams}
An \emph{embedded team}\index{embedded team} sits with one business (a desk or strategy group), reports into it and builds everything that business
needs. A \emph{central team}\index{central team} serves several businesses with one function, such as order gateways, market data or risk checks, and
reports into technology.
\end{definition}

The two designs trade the same things. \omterm{def:fm:engineering-organisation-design:teams}{Embedded teams} talk to their traders every day, change quickly and own their whole stack; they duplicate
what other desks also build, and each copy needs its own maintenance and on-call. \omterm{def:fm:engineering-organisation-design:teams}{Central teams} build one thing well for everyone and can
consolidate duplicates; every change a desk needs from them crosses a team boundary. Chapter 19's platform and \omterm{def:fm:technology-strategy:teams}{product teams} are the modern names
for the two poles, and Skelton and Pais's \emph{Team Topologies} (2019) their vocabulary; MacCormack, Rusnak and Baldwin (2012) tested Conway's
``mirroring'' of organisation in architecture empirically.

\begin{definition}[Coordination cost]\label{def:fm:engineering-organisation-design:coord}
The \emph{coordination cost}\index{coordination cost} of an organisation design is the work of keeping teams in step across the dependencies
between the services they own; the chapter measures it as the number of changes a week, to services that other teams' services depend on, which
those other teams must learn about, test against and schedule around.
\end{definition}

\section{The quant developer and other hybrid roles}

\begin{definition}[Quant developer]\label{def:fm:engineering-organisation-design:qd}
A \emph{quant developer}\index{quant developer} is an engineer who implements and runs trading models: who turns a researcher's prototype into
production code, owns its performance and correctness, and knows enough of the model to change it safely.
\end{definition}

Hybrid roles exist because the handoff between research and production is where models break. Book 12 described the machine-learning engineer
and the handoff specification between a model owner and the platform; the \omterm{def:fm:engineering-organisation-design:qd}{quant developer} is the same role for every kind of model. In an embedded
design they sit on the desk; in a central design they sit between the desk and the platform, and are the people the \omterm{def:fm:engineering-organisation-design:coord}{coordination cost} falls on.

\section{Ownership: every service has an owner}

\begin{definition}[Service owner, bus factor]\label{def:fm:engineering-organisation-design:owner}
A \emph{service owner}\index{service owner} is the named team, with a named lead, accountable for a service's correctness, changes, on-call and
retirement. The \emph{bus factor}\index{bus factor} of a service is the number of people who can operate and change it; the bus factor of an
organisation is the smallest over its services: the number of departures that would leave some service with no one who understands it.
\end{definition}

Ownership turns the design into data (\cref{lst:fm:engineering-organisation-design:measures}): a table of services with their owner, their
dependencies, how often each changes and how often it pages. From it three measures follow: the \omterm{def:fm:engineering-organisation-design:coord}{coordination cost}, the on-call load per engineer
and the \omterm{def:fm:engineering-organisation-design:owner}{bus factor}.

\omcode{code/firm/teamtopo/firm_teamtopo.py}{34}{61}{The three measures of an organisation design: cross-team dependencies weighted by change
frequency, pages per engineer by team, and the bus factor.}\label{lst:fm:engineering-organisation-design:measures}

\begin{proposition}[Designs move load, services create it]\label{prop:fm:engineering-organisation-design:load}
For a fixed set of services and engineers, the mean on-call load per engineer is the same in every design: total pages divided by engineers. A
design changes only how the load is distributed across teams, the \omterm{def:fm:engineering-organisation-design:coord}{coordination cost} and the \omterm{def:fm:engineering-organisation-design:owner}{bus factor}. Only removing or merging services, or
making them page less, lowers the mean.
\end{proposition}

\begin{proof}
Each service's pages are counted once, in its owner's team, whatever the owner; summing over teams gives the total pages, and dividing by the
engineers, who are all in some team, gives the mean.
\end{proof}

\section{Tutorial: three order gateways}\label{tut:fm:engineering-organisation-design:tut}

\begin{tutorial}
\textbf{Goal.} Compare an embedded and a central design for twenty services and fifteen engineers, then consolidate, then move one engineer and
see what breaks. \textbf{End state:} \cref{tab:fm:engineering-organisation-design:designs} and the load chart
(\cref{fig:fm:engineering-organisation-design:pages}).
\begin{tutsteps}
\item \textbf{The services.} Three desks (equities, futures, options), each with a strategy, an order gateway, a risk check, a market-data handler and
a pricing library; five shared services (security master, post-trade, monitoring, deployment, research platform); 33 dependencies; 34.5 changes and
6.2 pages a week in all (\texttt{fm\_teams}, illustrative).
\item \textbf{Embedded.} Each desk's team of four owns its five services; an infrastructure team of three owns the shared ones.
\item \textbf{Central.} Desk teams of two own the strategies and pricing; teams of two own the three gateways, the three risk checks and the three
market-data handlers; infrastructure as before.
\item \textbf{Consolidated.} The central design's natural next step: one gateway, one risk check and one market-data handler, owned by a trading
\omterm{def:fm:technology-strategy:teams}{platform team} of three; desk teams of three; fourteen services.
\item \textbf{The move.} One engineer leaves a desk (embedded) or the gateway team (central) for infrastructure.
\end{tutsteps}
\end{tutorial}

\begin{table}[ht]
\centering\small
\begin{tabular}{lrrrrrr}
\toprule
 & & & cross-team & coordination & max pages per & bus \\
design & teams & services & edges & (changes/week) & engineer-week & factor \\
\midrule
embedded & 4 & 20 & 18 & 22.5 & 0.43 & 3 \\
central & 7 & 20 & 30 & 33.0 & 0.90 & 2 \\
central, consolidated & 5 & 14 & 19 & 29.0 & 0.67 & 3 \\
\bottomrule
\end{tabular}
\caption{Three designs for fifteen engineers. Mean pages per engineer a week: 0.41 for both designs of the twenty services, 0.29 after
consolidation. Data: \texttt{fm\_teams.compare\_all}.}\label{tab:fm:engineering-organisation-design:designs}
\end{table}

The embedded design has the least coordination (22.5 changes a week crossing teams, on 18 edges), the most even on-call load (0.43 pages per
engineer a week at most) and a \omterm{def:fm:engineering-organisation-design:owner}{bus factor} of 3. Its cost is invisible in the table: it runs three of everything. Centralising the same twenty
services is the worst of both: 30 cross-team edges, 33 changes a week to coordinate, the market-data team carrying 0.90 pages per engineer a week, and
a \omterm{def:fm:engineering-organisation-design:owner}{bus factor} of 2, since every function now sits with two people. The mean load does not move
(\cref{prop:fm:engineering-organisation-design:load}): reorganising alone redistributes work and adds boundaries.

Consolidation is what the central design is for. Merging the three gateways, risk checks and market-data handlers leaves fourteen services and 4.3
pages a week, 31\% fewer; the mean load falls to 0.29 pages per engineer, the \omterm{def:fm:engineering-organisation-design:owner}{bus factor} returns to 3 with teams of three, and the desks gain an
engineer each. Coordination stays higher than in the embedded design (29.0 changes a week), because every desk now depends on one shared gateway:
that is the price of having one.

\begin{omfigure}
\begin{tikzpicture}
\begin{axis}[width=10.5cm,height=5.2cm,ybar,bar width=9pt,ylabel={\small count a week / edges},ymin=0,ymax=36,xtick={0,1,2},
  xticklabels={embedded,central,{central, consolidated}},xmin=-0.6,xmax=2.6,
  tick label style={font=\footnotesize},grid=major,grid style={gray!20},table/col sep=comma,
  legend cell align=left,legend style={font=\scriptsize,at={(0.5,-0.25)},anchor=north,/tikz/every even column/.append style={column sep=8pt}},legend columns=2]
\addplot[fill=omDef!55,draw=omDef] table[x=pos,y=coordination]{figdata/desk/21-engineering-organisation-design/designs.csv};
\addplot[fill=omProp!55,draw=omProp] table[x=pos,y=cross_edges]{figdata/desk/21-engineering-organisation-design/designs.csv};
\legend{coordination (changes a week),cross-team edges}
\end{axis}
\end{tikzpicture}

\omcaption{Coordination under the three designs: the number of dependencies that cross a team boundary, and the changes a week that cross them.
Centralising without consolidating adds boundaries; consolidating removes some of them again. Data: \texttt{fm\_teams.compare\_all}.}\label{fig:fm:engineering-organisation-design:coord}
\end{omfigure}

\section{On-call and its cost}

On-call load is measured in pages, but it is paid in hours. Google's site-reliability book caps on-call at a quarter of an engineer's time and
finds that an incident, with its root-cause analysis, remediation and follow-up, takes about six hours; at most two such incidents fit in a
twelve-hour shift (\cref{dat:fm:engineering-organisation-design:sre}). On the chapter's numbers the central design's market-data team spends
$0.9\times6=5.4$ hours a week each on incidents, 13.5\% of a forty-hour week; after one engineer leaves the gateway team, its remaining engineer
carries 1.5 pages a week, nine hours, 22.5\% of the week and close to the cap (\cref{fig:fm:engineering-organisation-design:pages}).

\begin{dated}{2026-09}{A published on-call budget}\label{dat:fm:engineering-organisation-design:sre}
Google's \emph{Site Reliability Engineering} (2016), chapter ``Being On-Call'': at least 50\% of site-reliability engineers' time goes to
engineering and no more than 25\% to on-call; handling an incident takes about six hours on average, so the maximum is two incidents per
twelve-hour shift. Book 15 applies these budgets to a trading platform's on-call rotation.
\end{dated}

\begin{omfigure}
\begin{tikzpicture}
\begin{axis}[width=10.5cm,height=5.4cm,xbar,bar width=6pt,xmin=0,xmax=1.1,xlabel={\small pages per engineer a week},
  symbolic y coords={embedded: desk,embedded: infrastructure,central: desk,central: gateways,central: risk,central: market data,central: infrastructure,consolidated: desk,consolidated: platform,consolidated: infrastructure},
  ytick=data,y dir=reverse,tick label style={font=\footnotesize},grid=major,grid style={gray!20},enlarge y limits=0.06,nodes near coords,
  nodes near coords style={font=\scriptsize,/pgf/number format/.cd,fixed,precision=2},point meta=x,table/col sep=comma]
\addplot[fill=omDef!55,draw=omDef] table[y=label,x=pages]{figdata/desk/21-engineering-organisation-design/pages.csv};
\end{axis}
\end{tikzpicture}

\omcaption{On-call load per engineer by team under the three designs (the three desk teams carry the same load in each design). The central design
concentrates pages on the small function teams. Data: \texttt{firm.teamtopo.pages}.}\label{fig:fm:engineering-organisation-design:pages}
\end{omfigure}

The move shows the \omterm{def:fm:engineering-organisation-design:owner}{bus factor} at work. In the embedded design, moving one engineer off the equities desk leaves its five services with three
people and raises that desk's load to 0.57 pages a week each; nothing is orphaned. In the central design, moving one engineer off the gateway team
leaves the three gateways with one person each: the \omterm{def:fm:engineering-organisation-design:owner}{bus factor} falls to 1, and one resignation would leave the firm unable to change how it sends
orders. In the consolidated design the same move leaves the trading platform with two people and a \omterm{def:fm:engineering-organisation-design:owner}{bus factor} of 2.

\begin{omfigure}
\begin{tikzpicture}[font=\footnotesize,>=stealth,t/.style={draw,rounded corners,fill=omDef!15,minimum width=1.6cm,align=center},
  s/.style={draw,rounded corners,fill=omThm!15,minimum width=1.5cm,align=center}]
\node[font=\small] at (-4.2,2.5) {embedded};
\node[t] (e) at (-6,1.5) {equities}; \node[t] (f) at (-4.2,1.5) {futures}; \node[t] (o) at (-2.4,1.5) {options};
\node[s] (ge) at (-6,0) {gateway 1}; \node[s] (gf) at (-4.2,0) {gateway 2}; \node[s] (go) at (-2.4,0) {gateway 3};
\draw[->] (e) -- (ge); \draw[->] (f) -- (gf); \draw[->] (o) -- (go);
\node[font=\small] at (3.6,2.5) {central, consolidated};
\node[t] (e2) at (1.8,1.5) {equities}; \node[t] (f2) at (3.6,1.5) {futures}; \node[t] (o2) at (5.4,1.5) {options};
\node[t,fill=omProp!15] (p) at (3.6,0) {trading platform team}; \node[s] (g) at (3.6,-1.3) {one gateway};
\draw[->,dashed] (e2) -- (p); \draw[->,dashed] (f2) -- (p); \draw[->,dashed] (o2) -- (p); \draw[->] (p) -- (g);
\draw[gray] (-0.3,-1.8) -- (-0.3,2.8);
\end{tikzpicture}

\omcaption{Conway's law in a trading firm: three desk teams build three gateways (left); a \omterm{def:fm:technology-strategy:teams}{platform team} builds one that all three desks
depend on (dashed, right). The second design has fewer services and more cross-team dependencies. Schematic.}\label{fig:fm:engineering-organisation-design:conway}
\end{omfigure}

\section{Measuring an engineering organisation}

The three measures are cheap to keep once ownership is data: \omterm{def:fm:engineering-organisation-design:coord}{coordination cost} from the dependency graph and change logs, on-call load from the
paging system, \omterm{def:fm:engineering-organisation-design:owner}{bus factor} from who has changed and operated each service in the last year. Forsgren, Humble and Kim's \emph{Accelerate} (2018)
proposes delivery measures (how often teams deploy, how long a change takes to reach production, how often changes fail, how fast service is
restored) that complete them.

\begin{method}[Designing the engineering organisation]\label{met:fm:engineering-organisation-design:design}
\begin{enumerate}
\item Keep the service catalogue as data: owner, dependencies, change frequency, page rate, who can operate it.
\item Draw teams around the communication the work needs (Conway's criterion): embed what changes with the desk, centralise what several desks
need in the same form, and consolidate what is duplicated.
\item Measure coordination, on-call load against a budget, and the \omterm{def:fm:engineering-organisation-design:owner}{bus factor}, before and after every reorganisation.
\item Keep every service's \omterm{def:fm:engineering-organisation-design:owner}{bus factor} at three or more; pair new people on the weakest services first.
\item Review the design when a desk or a venue is added: each is a new service or a new dependency.
\end{enumerate}
\end{method}

\section{Build: the organisation as a graph}\label{bld:fm:engineering-organisation-design:teamtopo}

\begin{build}
\bfield{Purpose} Services, dependencies, change and incident rates and team assignments as data, with the measures of a design.
\bfield{Interface} \texttt{firm.teamtopo}: \texttt{Service}, \texttt{Design}, \texttt{cross\_edges}, \texttt{coordination}, \texttt{pages},
\texttt{bus\_factor}, \texttt{move}, \texttt{report}.
\bfield{Rules} Every service has exactly one owning team; a service's \omterm{def:fm:engineering-organisation-design:owner}{bus factor} counts its team and named backups; the mean load is invariant
under reassignment.
\bfield{Acceptance tests} \texttt{code/firm/teamtopo/tests/}: cross-team edges and coordination on a three-service graph; pages by team; the \omterm{def:fm:engineering-organisation-design:owner}{bus
factor} before and after a move.
\bfield{Stretch} Knowledge measured from commit and deployment history; the cost of a reorganisation itself; optimising the assignment for a
weighted sum of the three measures.
\end{build}

\begin{omsources}
\item M.~E.~Conway, ``How do committees invent?'', \emph{Datamation}, April 1968.
\item A.~MacCormack, C.~Baldwin and J.~Rusnak, \emph{Research Policy} 41(8), 2012; M.~Skelton and M.~Pais, \emph{Team Topologies}, 2019.
\item B.~Beyer et al. (eds), \emph{Site Reliability Engineering}, 2016; N.~Forsgren, J.~Humble and G.~Kim, \emph{Accelerate}, 2018.
\end{omsources}

\section{Exercises}

\begin{exercise}[$\star$]\label{exo:fm:engineering-organisation-design:1}
State Conway's thesis and his criterion for organising a design effort.
\end{exercise}

\begin{exercise}[$\star$]\label{exo:fm:engineering-organisation-design:2}
What is the mean on-call load per engineer in each design, and why is it the same for the embedded and central designs?
\end{exercise}

\begin{exercise}[$\star$]\label{exo:fm:engineering-organisation-design:3}
How many hours a week does the central market-data team spend on incidents, at six hours an incident, and what share of the on-call cap is it?
\end{exercise}

\begin{exercise}[$\star\star$]\label{exo:fm:engineering-organisation-design:4}
Why does centralising the same twenty services raise coordination from 22.5 to 33.0 changes a week?
\end{exercise}

\begin{exercise}[$\star\star$]\label{exo:fm:engineering-organisation-design:5}
What happens to the \omterm{def:fm:engineering-organisation-design:owner}{bus factor} in each design when one engineer moves, and which services are at risk?
\end{exercise}

\begin{exercise}[$\star\star$]\label{exo:fm:engineering-organisation-design:6}
Where should the \omterm{def:fm:engineering-organisation-design:qd}{quant developers} sit in each design, and why?
\end{exercise}

\begin{exercise}[$\star\star\star$]\label{exo:fm:engineering-organisation-design:7}
\emph{Coding.} In the consolidated design, move one engineer from the trading platform to infrastructure. What are the \omterm{def:fm:engineering-organisation-design:owner}{bus factor} and the largest
load?
\end{exercise}

\begin{exercise}[$\star\star\star$]\label{exo:fm:engineering-organisation-design:8}
\emph{Find the flaw.} ``Centralising our engineers will cut our on-call load.''
\end{exercise}

\section{Problem: Three Order Gateways}

\begin{problem}[{Weekend problem --- three order gateways}]\label{pb:fm:engineering-organisation-design:1}
A new chief technology officer inherits three desks with their own engineers and three of everything, and must propose an organisation.

\medskip\noindent\textbf{Part I --- The principles.}
\begin{enumerate}
\item State Conway's thesis.
\item Define embedded and \omterm{def:fm:engineering-organisation-design:teams}{central teams} and their trade-off.
\item Define a \omterm{def:fm:engineering-organisation-design:qd}{quant developer} and say why the role exists.
\item Define \omterm{def:fm:engineering-organisation-design:coord}{coordination cost}, \omterm{def:fm:engineering-organisation-design:owner}{service owner} and \omterm{def:fm:engineering-organisation-design:owner}{bus factor}.
\end{enumerate}

\medskip\noindent\textbf{Part II --- The model.}
\begin{enumerate}[resume]
\item Describe the services and their dependencies.
\item State and prove \cref{prop:fm:engineering-organisation-design:load}.
\item Give the three measures for the embedded and central designs.
\item Give them for the consolidated design and explain what consolidation buys and costs.
\end{enumerate}

\medskip\noindent\textbf{Part III --- Stress.}
\begin{enumerate}[resume]
\item What is the published on-call budget, and how close does each design come to it?
\item Move one engineer in each design and give the result.
\item Which services should be paired first, and why?
\item How would you measure the \omterm{def:fm:engineering-organisation-design:owner}{bus factor} from data?
\end{enumerate}

\medskip\noindent\textbf{Part IV --- The proposal.}
\begin{enumerate}[resume]
\item Which services should stay embedded, which be centralised, which be consolidated?
\item How would you sequence the change?
\item What delivery measures would you add?
\item What is the risk of a single shared gateway, and how would you manage it?
\item How often should the design be reviewed?
\item Why is reorganising without consolidating the worst option in the model?
\item State the \emph{named result}: \omterm{def:fm:engineering-organisation-design:coord}{coordination cost} and pages per engineer per week under the embedded and the central design, and the \omterm{def:fm:engineering-organisation-design:owner}{bus factor}
of each.
\item In two sentences, write the proposal.
\end{enumerate}
\end{problem}

\section{Interview questions}

\begin{interviewq}[$\star$ \iqroles{developer}]\label{iq:fm:engineering-organisation-design:1}
What is Conway's law, and where have you seen it?
\end{interviewq}

\begin{interviewq}[$\star$ \iqroles{developer, researcher}]\label{iq:fm:engineering-organisation-design:2}
What does a \omterm{def:fm:engineering-organisation-design:qd}{quant developer} do that a researcher and a software engineer do not?
\end{interviewq}

\begin{interviewq}[$\star\star$ \iqroles{developer}]\label{iq:fm:engineering-organisation-design:3}
How would you compute the \omterm{def:fm:engineering-organisation-design:owner}{bus factor} of a codebase?
\end{interviewq}

\begin{interviewq}[$\star\star$ \iqroles{developer}]\label{iq:fm:engineering-organisation-design:4}
Your team gets two pages a night. What do you do?
\end{interviewq}

\begin{interviewq}[$\star\star$ \iqroles{trader, developer}]\label{iq:fm:engineering-organisation-design:5}
Should each trading desk have its own engineers?
\end{interviewq}

\begin{interviewq}[$\star\star\star$ \iqroles{developer, researcher}]\label{iq:fm:engineering-organisation-design:6}
Formulate the assignment of services to teams as an optimisation problem.
\end{interviewq}
